\documentclass[../../../include/open-logic-section]{subfiles}

\begin{document}

\olfileid{sth}{spine}{foundation}

\olsection{基础公理}

我们只算\emph{几乎}完成——而非\emph{彻底}完成——因为 $\ZFminus$ 中没有任何内容能保证\emph{每个}集合都在某个 $V_\alpha$ 中，即每个集合都在某个阶段形成。

现在，从（数学上）相当直接的意义来看，我们并不\emph{在乎}是否存在层级之外的集合。（如果真有这种集合，我们直接忽略它们即可。）但我们正是通过“每个集合都在某个阶段形成”这一想法来确立我们的集合\emph{概念}的（参见 \olref[z][story]{sec} 中的 \stageshier{}）。因此，我们希望排除集合落在层级之外的可能性。相应地，我们必须添加一条新公理，以确保每个集合都出现在层级中的某处。

既然 $V_\alpha$ 就是我们的阶段，我们或许会考虑直接将以下内容作为公理加入：

\begin{defish}
\emph{正则性公理。} $\forall A \exists \alpha\, A \subseteq V_\alpha$
\end{defish}

这将是一种完全合理的做法。然而，由于下一节将说明的原因，我们将转而采用另一条公理：

\begin{axiom}[基础公理]
$(\forall A \neq \emptyset)(\exists B \in A)A \cap B = \emptyset$。
\end{axiom}

经过一番努力，我们可以（在 $\ZFminus$ 中）证明基础公理蕴涵正则性公理：
\begin{defn}
对每个集合 $A$，令：
\begin{align*}
	\text{cl}_0(A) &= A,\\
	\text{cl}_{n+1}(A) &= \bigcup \text{cl}_n(A),\\
	\text{trcl}(A) &= \bigcup_{n < \omega} \text{cl}_{n}(A).
\end{align*}
我们称 $\text{trcl}(A)$ 为 $A$ 的\emph{传递闭包}。
\end{defn}
\noindent
“传递闭包”这一名称是恰当的：
\begin{prop}\ollabel{subsetoftrcl}
$A \subseteq \trcl{A}$，并且 $\trcl{A}$ 是一个传递集。
\end{prop}

\begin{proof}
显然 $A = \text{cl}_0(A) \subseteq \text{trcl}(A)$。并且，如果 $x \in b \in \trcl{A}$，那么对某个 $n$ 有 $b \in \text{cl}_n(A)$，因此 $x
\in \text{cl}_{n+1}(A) \subseteq \text{trcl}(A)$。
\end{proof}

\begin{lem}\ollabel{lem:TransitiveWellFounded}
如果 $A$ 是一个传递集，则存在某个 $\alpha$ 使得 $A \subseteq V_\alpha$。
\end{lem}

\begin{proof}
回顾 \olref[ordinals][opps]{defsupstrict} 中“$\supstrict(X)$”的定义，定义两个集合：
\begin{align*}
	D  &= \Setabs{x \in A}{\forall \delta\ x \nsubseteq V_\delta}\\
	\alpha &= \supstrict\Setabs{\delta}{(\exists x \in A)
	(x \subseteq V_\delta \land (\forall \gamma \in \delta)x \nsubseteq V_\gamma)}
\end{align*}
假设 $D = \emptyset$。那么，若 $x \in A$，则存在某个 $\delta$ 使得 $x \subseteq V_\delta$，并且根据序数的良序性，$(\forall \gamma \in \delta)x \nsubseteq V_\gamma$；因此 $\delta \in \alpha$，从而由 \olref[spine][Valphabasic]{Valphabasicprops} 知 $x \in V_\alpha$。于是 $A \subseteq V_{\alpha}$ 得证。

因此，只需证明 $D = \emptyset$。用归谬法，假设不然。由基础公理，存在某个 $B \in D \subseteq A$ 使得 $D \cap B = \emptyset$。若 $x \in B$，则由于 $A$ 是传递的，有 $x \in A$；又因 $x \notin D$，可知 $\exists \delta\ x \subseteq V_\delta$。现在令
\[
	\beta = \supstrict\Setabs{\delta}{(\exists x \in b)(x \subseteq V_\delta \land (\forall \gamma < \delta)x \nsubseteq V_\gamma)}.
\]
同前，$B \subseteq V_\beta$，这与 $B \in D$ 矛盾。
\end{proof}

\begin{thm}\ollabel{zfentailsregularity}
正则性公理成立。
\end{thm}

\begin{proof}
固定 $A$；由 \olref{subsetoftrcl} 知 $A \subseteq \trcl{A}$，且 $\trcl{A}$ 是传递集。于是根据 \olref{lem:TransitiveWellFounded}，存在某个 $\alpha$ 使得 $A \subseteq \trcl{A} \subseteq V_\alpha$。
\end{proof}

这些结果表明，$\ZFminus$ 可以证明条件式 $\emph{基础公理}\Rightarrow\emph{正则性公理}$。在 \olref[spine][rank]{zfminusregularityfoundation} 中，我们将证明 $\ZFminus$ 可以证明 $\emph{正则性公理}\Rightarrow\emph{基础公理}$。这样一来，在 $\ZFminus$ 下，基础公理与正则性公理是\emph{等价}的。但这就意味着，给定 $\ZFminus$，我们可以通过指出基础公理等价于正则性公理来证成基础公理。而我们又能依据 \stageshier{} 直接证成正则性公理。

\end{document}